% TM.tex (RMCG20260121)

\input RCstyle

\input doubleco
\darkcolors

% Every point a period
\count255=`A
\loop
\sfcode\count255=1000
\ifnum\count255<`Z\advance\count255 1
\repeat


% define#1|#2|#3|#4|
%    #1 if empty write {\em #2}, else (if is *) do not write
%    #2 a unique label - to be ordered alphabetically - preferred ASCII
%    #3 the name of the concept typographically, empty = #1
%    #4 its definition goes to macro \^^06#1

\catcode`\^^06=12
\def\define#1|#2|#3|#4|{\def\1{#1}\def\3{#3}\ifx\3\empty \def\3{#2}\fi
 \expandafter\gdef\csname ^^06#2\endcsname{#4}%
 \lbl{#2}{\3}\ndx{#2}{D}\ifx\1\empty {\em\3}\fi}

 \def\TMndxline#1#2#3#4#5{\def\2{#2}%
  \ifx\2\D\par \hangindent20pt\noindent
   {\em\ref{#1}}\quad(\refsc{#1})\quad\ignorespaces
  \csname^^06#1\endcsname\ignorespaces\else\ignorespaces\fi}

\def\definitions{\par\begingroup
 \everypar{}\parno=0\def\secid{}\noindent\dest
 {\fonttwo Definitions\toc1{Definitions}\lbl{Definitions}{Definitions}}
 \par\def\D{D}\let\ndxline=\TMndxline
 \input auxiliar.abc
 \par \endgroup}


\chardef\diffchar=`!

\catcode`\@=11

\newcount\sssecno \sssecno=0

\outer\def\subsubsection#1{\vskip0pt plus 60pt\penalty-250\vskip0pt plus-60pt
 \ifnum\parno=0 \vskip 0.5pc plus 3pt minus 3pt\else\vskip 1pc plus 6pt minus 6pt\fi
 \ifx\n@xtlbl\relax \def\1{#1}\else\let\1=\n@xtlbl \let\n@xtlbl=\relax \fi
 \everypar{}\advance\sssecno1 \thno=0 \parno=0
 \def\secid{\presec\the\secno.\the\ssecno.\the\sssecno}\noindent\dest
 {\setbox2=\hbox{\sl\secid\enspace}\hangindent\wd2
  \rightskip=0pt plus 6em
 \sl\secid\enspace#1\toc3{#1}\lbl{\1}{#1}\par}
 \everypar{\numberedpars}}

\def\file{\ifvmode\vskip\parskip\else\par\fi
 \begingroup \lineskip=3pt \ifdim\parindent<20pt \parindent=20pt \fi
 \catcode`\~=12 \catcode`\"=12 \catcode`\_=12 \catcode`\\=12 \@file}
\def\@file#1{\index{#1}1
 \hbox{\indent\vbox{\hrule\hbox{\strut\color{Blue}\tt#1\endcolor}\hrule}}\nobreak
 \let\verbatimoptions=\tenttlistingoptions \setupverbatim \input#1
 \dimen0=\prevdepth \advance\dimen0by-3pt \prevdepth=\dimen0
 \everypar{}\noindent\endgroup \ignorespaces}
\def\tenttlistingoptions{\tentt
 \baselineskip=12pt plus0.1pt minus0.1pt
 \parindent=40pt \rightskip=0pt plus 3cm
 \everypar{\hangindent60pt \llap{\sevenrm\the\inputlineno\quad}}}

\def\codeoptions{\tentt \baselineskip=12pt plus0.1pt minus0.1pt
 \parindent=20pt}
\def\begincode{\ifvmode \vskip\parskip \else \par\smallskip\fi
  \begingroup \catcode`\~=12 \catcode`\"=12 \catcode`\_=12 \catcode`\\=12
  \let\verbatimoptions\codeoptions\setupverbatim\d@code}
\catcode`/=0 \catcode`\\=12
 /def/d@code{/def/next##1##2\endcode{##2/everypar{}/smallskip%
   /noindent/endgroup/ignorespaces}/next}
/catcode`\=0 /catcode`\/=12

% \plisting#1
% #2
% #3
% prints file #1 contents verbatim
% from the first line equal to #2
% to the next line equal to #3
\def\plisting#1{\ifvmode\vskip\parskip\else\par\fi \begingroup
 \def~{\string~}\def\1{#1}\beforeplisting
 \def\d@next{\expandafter\pa@ss\input #1 }\setupverbatim\getm@rks}
{\obeylines\gdef\getm@rks^^M#1^^M#2^^M{%
  \gdef\beginm@rk{#1}\gdef\endm@rk{#2}\d@next}}
{\obeylines\gdef\pa@ss#1^^M{\def\next{#1}\ifx\next\beginm@rk %
  #1\par \let\next\pr@nt \else \let\next\pa@ss\fi \next}}
\def\endpr@nt{\afterplisting\endgroup\endinput}
{\obeylines\gdef\pr@nt#1^^M{\def\next{#1}\ifx\next\endm@rk %
  \let\next\endpr@nt \else #1\par \let\next\pr@nt \fi \next}}
\let\beforeplisting=\relax \let\afterplisting=\relax

{\obeylines\gdef\getm@rk#1^^M{\def\endm@rk{#1}\d@next}}

\def\pfile{\begingroup\let\beforeplisting=\tenttplisting
 \let\afterplisting=\endgroup\plisting}
\def\tenttplisting{\ifdim\lastskip<\smallskipamount
 \removelastskip \smallskip\fi \index{\1}1
 \let\verbatimoptions=\tenttlistingoptions
 \ifx\1\@pfile \else \xdef\@pfile{\1}%
  \hbox{\indent\vbox{\hrule\hbox{\strut\color{Blue}\tt\1\endcolor}\hrule}}\nobreak
 \fi}

\catcode`\@=12

\newread\situations
\newcount\sitno

\def\nextsituation{\read\situations to\situation \situation}
\def\skipsituation#1{\sitno=#1\loop \read\situations to\situation
 \advance\sitno-1 \ifnum\sitno>1\repeat}


\newread\lines
\newcount\linno

\def\listline#1#2{\openin\lines=#1
 \linno=0\loop \read\lines to\rline
 \advance\linno1 \ifnum\linno<#2\repeat
 \hbox{\kern40pt\llap{\sevenrm #2\quad}\tentt\rline}%
 \closein\lines}


\def\M#1{\hbox{\tt#1}}
\def\T#1{\\\null\kern\oldparindent\M{#1}}

\def\Table#1{\vtop{%
\halign{\strut\vrule\enspace\M{##}\hfil\enspace
 &\enspace\hfil\M{##}\hfil\enspace
 &\vrule\enspace\hfil\M{##}\hfil\enspace
 &\enspace\hfil\M{##}\hfil\enspace
 &\enspace\hfil\M{##}\hfil\enspace\vrule\cr
 \noalign{\vskip-11pt\hrule}
 {\sc c}&{\sc r}&{\sc n}&{\sc w}&{\sc m}\cr
 \noalign{\hrule}\noalign{\kern\jot}\noalign{\hrule}
 #1\crcr
 \noalign{\hrule}
}}}

\def\tape#1#2{\line{\qquad\tt #1\hfil\tentt #2}}

\def\TM#1<#2>{\def\1{#1}\ifx\empty\1\else{\cal #1}\fi
 \langle{\frak #2}\rangle}
\def\RS#1<#2>{\def\1{#1}\ifx\empty\1\else{\cal #1}\fi
 \langle\hbox{\tt #2}\rangle}

\def\label#1{\leavevmode\dest\lbl{#1}{label}\ignorespaces}

\MTfigures

\def\MTbeginfig{\begingroup \catcode`\"=12 \MTbeginchar}
\def\MTendfig;{\MTendchar; \endgroup}

\MTcode
def state(suffix s)(expr side) =
 x.s.iii = x.s + side;   y.s.iii = y.s;
 x.s.ix = x.s - side;    y.s.ix = y.s;
 x.s.i = x.s + side/2;   y.s.i = y.s + side * sqrt(3)/2;
 x.s.v = x.s + side/2;   y.s.v = y.s - side * sqrt(3)/2;
 x.s.vii = x.s - side/2; y.s.vii = y.s - side * sqrt(3)/2;
 x.s.xi = x.s - side/2;  y.s.xi = y.s + side * sqrt(3)/2;
 z.s.ii = 1/2[z.s.i,z.s.iii];
 z.s.iv = 1/2[z.s.iii,z.s.v];
 z.s.vi = 1/2[z.s.v,z.s.vii];
 z.s.viii = 1/2[z.s.vii,z.s.ix];
 z.s.x = 1/2[z.s.ix,z.s.xi];
 z.s.xii = 1/2[z.s.xi,z.s.i];
 draw (z.s.i) -- (z.s.iii) -- (z.s.v) -- (z.s.vii) -- (z.s.ix) -- (z.s.xi) -- cycle;
enddef;

\MTcode
def circle(suffix s)(expr diameter) =
 x.s.l = x.s - diameter/2;  x.s.r = x.s + diameter/2;
 y.s.t = y.s + diameter/2;  y.s.b = y.s - diameter/2;
 erase fill (fullcircle scaled (x.s.r - x.s.l) shifted z.s);
 draw (fullcircle scaled (x.s.r - x.s.l) shifted z.s);
enddef;

\MTcode
defaultfont:="cmtt12";
pickup pencircle scaled 0.4pt; thinpen:=savepen;
pickup pencircle scaled 0.8pt; medpen:=savepen;
pickup pencircle scaled 1.2pt; thickpen:=savepen;

\def\MTresetpa;{\MTline{ path pa[],
 pa[]O,pa[]I,pa[]A,pa[]B,pa[]X,pa[]Y,pa[]M,pa[]N,pa[]L,pa[]R,pa[]S,pa[]T,pa[]Z;}}

\def\MTstate(#1)"#2""#3";{%
  \MTline{ state(#1,6mm);}
  \MTline{ z#1.lbl1 = z#1+(0,3pt);}
  \MTline{ label("#2", z#1.lbl1);}
  \MTline{ z#1.lbl2 = z#1.vi+(0,8pt);}
  \MTline{ label("#3", z#1.lbl2);}
}

\def\MTinput(#1.#2)(#3.#4)"#5""#6";{{%
  \def\0{0}\def\1{1}\def\9{*}\def\5{#5}%
  \def\6{#6}\ifx\5\6\def\6{}\fi
  \ifx\5\0\def\5{O}\fi \ifx\5\1\def\5{I}\fi \ifx\5\9\def\5{Z}\fi
  \MTline{ z#1.\5.in = 1.2[z#1,z#1.#2];}
  \MTline{ z#1.\5.lbl1 = z#1.\5.in;}
  \MTline{ z#1.\5.o = z#1.#2;}
  \MTline{ z#1.\5.d = z#3.#4;}
  \MTline{ pa#1.\5 = z#1.\5.o}
  \MTline{   if known z#1.\5.p: .. z#1.\5.p fi}
  \MTline{   if known z#1.\5.q: .. z#1.\5.q fi}
  \MTline{   if known z#1.\5.r: .. z#1.\5.r fi}
  \MTline{   if known z#1.\5.t: .. z#1.\5.t fi}
  \MTline{   .. z#1.\5.d;}
  \MTline{ drawarrow pa#1.\5;}
  \MTline{ circle(#1.\5.in)(4mm);}
  \MTline{ label("#5", z#1.\5.lbl1);}
  \ifx\6\empty\else
  \MTline{ z#1.\5.lbl2 = point (arctime 15pt of pa#1.\5) of pa#1.\5;}
  \MTline{ erase fill (fullcircle scaled 10pt shifted z#1.\5.lbl2);}
  \MTline{ label("#6", z#1.\5.lbl2);}
  \fi
}}

\def\MTauto(#1)(#2,#3)"#4""#5";{{%
  \def\0{0}\def\1{1}\def\9{*}\def\4{#4}%
  \ifx\4\0\def\4{O}\fi \ifx\4\1\def\4{I}\fi \ifx\4\9\def\4{Z}\fi
  \MTline{ z#1.\4.in = 1.2[z#1,z#1.#2];}
  \MTline{ z#1.\4.lbl1 = z#1.\4.in;}
  \MTline{ z#1.\4.o = z#1.#2;}
  \MTline{ z#1.\4.d = z#1.#3;}
  \MTline{ z#1.\4.p = 2[z#1,z#1.\4.o];}
  \MTline{ z#1.\4.q = 2[z#1,z#1.\4.d];}
  \MTline{ pa#1.\4 = z#1.\4.o{z#1.#2-z#1} .. z#1.\4.p .. z#1.\4.q .. z#1.\4.d{z#1-z#1.#3};}
  \MTline{ z#1.\4.lbl2 = point (arctime 15pt of pa#1.\4) of pa#1.\4;}
  \MTline{ drawarrow pa#1.\4;}
  \MTline{ circle(#1.\4.in)(4mm);}
  \MTline{ label("#4", z#1.\4.lbl1);}
  \MTline{ erase fill (fullcircle scaled 10pt shifted z#1.\4.lbl2);}
  \MTline{ label("#5", z#1.\4.lbl2);}
}}

\catcode`\"=13 \def"{\verb"}

\endinput
